\documentclass[10pt,a4paper]{article}
\usepackage{amsmath,amssymb,amsthm,graphicx}
% Independent article wrapper; no publisher class or third-party style redistribution.
\usepackage[margin=25mm]{geometry}
\usepackage{fontspec}
\setmainfont{Latin Modern Roman}
\usepackage{amsmath,amssymb,amsthm,mathtools,graphicx}
\usepackage{hyperref}
\hypersetup{hidelinks}
\newtheorem{cdrthm}{Theorem}[section]
\newtheorem{theo}[cdrthm]{Theorem}
\newtheorem{prop}[cdrthm]{Proposition}
\newtheorem{lemm}[cdrthm]{Lemma}
\newtheorem{lemma}[cdrthm]{Lemma}
\newtheorem{coro}[cdrthm]{Corollary}
\theoremstyle{definition}
\newtheorem{defi}[cdrthm]{Definition}
\theoremstyle{remark}
\newtheorem{rema}[cdrthm]{Remark}
% Original author macros, published lines 70--437 and 1563--1564.
\theoremstyle{plain}
\newtheorem{problem}{Problem}
\newtheorem{fact}{Fact}

\newcommand{\psp}[1]{{{\bf P}^{#1}}}
\newcommand{\psr}[1]{{\bf P}(#1)}

\newcommand{\ini}[1]{\text{in}(#1)}
\newcommand{\gin}[1]{\text{gin}(#1)}
\newcommand{\kr}{{\Bbbk}}
\newcommand{\pd}{\partial}
\newcommand{\vardel}{\partial}
\renewcommand{\tt}{{\bf t}}



\newcommand{\coh}{{{\text{{\rm coh}}}}}

\newcommand{\modv}[1]{{#1}\text{-{mod}}}
\newcommand{\modstab}[1]{{#1}-\underline{\text{mod}}}


\newcommand{\sut}{{}^{\tau}}
\newcommand{\sumit}{{}^{-\tau}}
\newcommand{\til}{\thicksim}


\newcommand{\totp}{\text{Tot}^{\prod}}
\newcommand{\dsum}{\bigoplus}
\newcommand{\dprod}{\prod}
\newcommand{\lsum}{\oplus}
\newcommand{\lprod}{\Pi}

\newcommand{\La}{{\Lambda}}

\newcommand{\sirstj}{\circledast}

\newcommand{\she}{\EuScript{S}\text{h}}
\newcommand{\cm}{\EuScript{CM}}
\newcommand{\cmd}{\EuScript{CM}^\dagger}
\newcommand{\cmri}{\EuScript{CM}^\circ}
\newcommand{\cler}{\EuScript{CL}}
\newcommand{\clerd}{\EuScript{CL}^\dagger}
\newcommand{\clerri}{\EuScript{CL}^\circ}
\newcommand{\gor}{\EuScript{G}}
\newcommand{\cF}{\mathcal{F}}
\newcommand{\cG}{\mathcal{G}}

\newcommand{\cM}{\mathcal{M}}
\newcommand{\cE}{\mathcal{E}}

\newcommand{\cI}{\mathcal{I}}
\newcommand{\cP}{\mathcal{P}}
\newcommand{\cK}{\mathcal{K}}

\newcommand{\cS}{\mathcal{S}}
\newcommand{\cC}{\mathcal{C}}
\newcommand{\cO}{\mathcal{O}}
\newcommand{\cJ}{\mathcal{J}}
\newcommand{\cU}{\mathcal{U}}

\newcommand{\cQ}{\mathcal{Q}}
\newcommand{\cX}{\mathcal{X}}
\newcommand{\cY}{\mathcal{Y}}
\newcommand{\cZ}{\mathcal{Z}}
\newcommand{\cV}{\mathcal{V}}

\newcommand{\mm}{\mathfrak{m}}

\newcommand{\dlim} {\varinjlim}
\newcommand{\ilim} {\varprojlim}


\newcommand{\CM}{\text{CM}}
\newcommand{\Mon}{\text{Mon}}

\newcommand{\Kom}{\text{Kom}}

\newcommand{\EH}{{\mathbf H}}
\newcommand{\res}{\text{res}}
\newcommand{\Hom}{{\rm{Hom}}}
\newcommand{\inhom}{{\underline{\text{Hom}}}}
\newcommand{\Ext}{\text{Ext}}
\newcommand{\Tor}{\text{Tor}}
\newcommand{\ghom}{\mathcal{H}om}
\newcommand{\gext}{\mathcal{E}xt}
\newcommand{\id}{\text{{id}}}
\newcommand{\im}{\text{im}\,}
\newcommand{\codim} {\text{codim}\,}
\newcommand{\resol}{\text{resol}\,}
\newcommand{\rank}{\text{rank}\,}
\newcommand{\lpd}{\text{lpd}\,}
\newcommand{\coker}{\text{coker}\,}
\newcommand{\supp}{\text{supp}\,}
\newcommand{\Ad}{A_\cdot}
\newcommand{\Bd}{B_\cdot}
\newcommand{\Fd}{F_\cdot}
\newcommand{\Gd}{G_\cdot}

\newcommand{\sus}{\subseteq}
\newcommand{\sups}{\supseteq}
\newcommand{\pil}{\rightarrow}
\newcommand{\vpil}{\leftarrow}
\newcommand{\rpil}{\leftarrow}
\newcommand{\lpil}{\longrightarrow}
\newcommand{\inpil}{\hookrightarrow}
\newcommand{\pils}{\twoheadrightarrow}
\newcommand{\projpil}{\dashrightarrow}
\newcommand{\dotpil}{\dashrightarrow}
\newcommand{\adj}[2]{\overset{#1}{\underset{#2}{\rightleftarrows}}}
\newcommand{\mto}[1]{\stackrel{#1}\longrightarrow}
\newcommand{\vmto}[1]{\stackrel{#1}\longleftarrow}
\newcommand{\mtoelm}[1]{\stackrel{#1}\mapsto}
\newcommand{\bihom}[2]{\overset{#1}{\underset{#2}{\rightleftarrows}}}
\newcommand{\eqv}{\Leftrightarrow}
\newcommand{\impl}{\Rightarrow}

\newcommand{\iso}{\cong}
\newcommand{\te}{\otimes}
\newcommand{\into}[1]{\hookrightarrow{#1}}
\newcommand{\ekv}{\Leftrightarrow}
\newcommand{\equi}{\simeq}
\newcommand{\isopil}{\overset{\cong}{\lpil}}
\newcommand{\equipil}{\overset{\equi}{\lpil}}
\newcommand{\ispil}{\isopil}
\newcommand{\vvi}{\langle}
\newcommand{\hvi}{\rangle}
\newcommand{\susneq}{\subsetneq}
\newcommand{\sgn}{\text{sign}}


\newcommand{\xd}{\check{x}}
\newcommand{\ortog}{\bot}
\newcommand{\tL}{\tilde{L}}
\newcommand{\tM}{\tilde{M}}
\newcommand{\tH}{\tilde{H}}
\newcommand{\tvH}{\widetilde{H}}
\newcommand{\tvh}{\widetilde{h}}
\newcommand{\tV}{\tilde{V}}
\newcommand{\tS}{\tilde{S}}
\newcommand{\tT}{\tilde{T}}
\newcommand{\tR}{\tilde{R}}
\newcommand{\tf}{\tilde{f}}
\newcommand{\ts}{\tilde{s}}
\newcommand{\tp}{\tilde{p}}
\newcommand{\tr}{\tilde{r}}
\newcommand{\tfst}{\tilde{f}_*}
\newcommand{\empt}{\emptyset}
\newcommand{\bfa}{{\mathbf a}}
\newcommand{\bfb}{{\mathbf b}}
\newcommand{\bfd}{{\mathbf d}}
\newcommand{\bfl}{{\mathbf \ell}}
\newcommand{\bfx}{{\mathbf x}}
\newcommand{\bfm}{{\mathbf m}}
\newcommand{\bfv}{{\mathbf v}}
\newcommand{\bft}{{\mathbf t}}
\newcommand{\bbfa}{{\mathbf a}^\prime}
\newcommand{\la}{\lambda}
\newcommand{\bfen}{{\mathbf 1}}
\newcommand{\bfe}{{\mathbf 1}}
\newcommand{\ep}{\epsilon}
\newcommand{\en}{r}
\newcommand{\tu}{s}
\newcommand{\Sym}{\text{Sym}}

\newcommand{\ome}{\omega_E}

\newcommand{\bevis}{{\bf Proof. }}
\newcommand{\demofin}{\qed \vskip 3.5mm}
\newcommand{\nyp}[1]{\noindent {\bf (#1)}}
\newcommand{\demo}{{\it Proof. }}
\newcommand{\demodone}{\demofin}
\newcommand{\parg}{{\vskip 2mm \addtocounter{theorem}{1}  
                   \noindent {\bf \thetheorem .} \hskip 1.5mm }}

\newcommand{\lcm}{{\text{lcm}}}

\newcommand{\dl}{\Delta}
\newcommand{\cdel}{{C\Delta}}
\newcommand{\cdelp}{{C\Delta^{\prime}}}
\newcommand{\dlst}{\Delta^*}
\newcommand{\Sdl}{{\mathcal S}_{\dl}}
\newcommand{\lk}{\text{lk}}
\newcommand{\lkd}{\lk_\Delta}
\newcommand{\lkp}[2]{\lk_{#1} {#2}}
\newcommand{\del}{\Delta}
\newcommand{\delr}{\Delta_{-R}}
\newcommand{\dd}{{\dim \del}}
\newcommand{\Del}{\Delta}

\newcommand{\bb}{{\mathbf b}}
\newcommand{\cc}{{\mathbf c}}
\newcommand{\xx}{{\mathbf x}}
\newcommand{\yy}{{\mathbf y}}
\newcommand{\zz}{{\mathbf z}}
\newcommand{\mv}{{\xx^{\aa_v}}}
\newcommand{\mF}{{\xx^{\aa_F}}}

\newcommand{\Symm}{\text{Sym}}
\newcommand{\pnm}{{\mathbf P}^{n-1}}
\newcommand{\opnm}{{\go_{\pnm}}}
\newcommand{\ompnm}{\omega_{\pnm}}

\newcommand{\pn}{{\mathbf P}^n}
\newcommand{\hele}{{\mathbb Z}}
\newcommand{\nat}{{\mathbb N}}
\newcommand{\rasj}{{\mathbb Q}}
\newcommand{\bfone}{{\mathbf 1}}

\newcommand{\dt}{\bullet}

\newcommand{\disk}{\scriptscriptstyle{\bullet}}

\newcommand{\cxF}{F_\dt}
\newcommand{\pol}{f}

\newcommand{\Rn}{{\mathbb R}^n}
\newcommand{\An}{{\mathbb A}^n}
\newcommand{\frg}{\mathfrak{g}}
\newcommand{\PW}{{\mathbb P}(W)}

\newcommand{\pos}{{\mathcal Pos}}
\newcommand{\g}{{\gamma}}

\newcommand{\Vaa}{V_0}
\newcommand{\Bp}{B^\prime}
\newcommand{\Bpp}{B^{\prime \prime}}
\newcommand{\bbp}{\mathbf{b}^\prime}
\newcommand{\bbpp}{\mathbf{b}^{\prime \prime}}
\newcommand{\bp}{{b}^\prime}
\newcommand{\bpp}{{b}^{\prime \prime}}

\newcommand{\oLa}{\overline{\Lambda}}
\newcommand{\ov}[1]{\overline{#1}}
\newcommand{\ovv}[1]{\overline{\overline{#1}}}
\newcommand{\tm}{\tilde{m}}
\newcommand{\po}{\bullet}

\newcommand{\surj}[1]{\overset{#1}{\twoheadrightarrow}}
\newcommand{\Supp}{\text{Supp}}

\newcommand{\ZZ}{{\mathbb Z}}
\newcommand{\NN}{{\mathbb N}}
\newcommand{\RR}{{\mathbb R}}

\newcommand{\oR}{\overline{R}}

\newcommand{\bfu}{{\mathbf u}}
\newcommand{\nn}{{\mathbf n}}
\newcommand{\oa}{\overline{a}}
\newcommand{\cop}{\text{cop}}

\newcommand{\opos}{\text{op}\,}

\newcommand{\ngmi}{\text{neg}}
\newcommand{\up}{\text{up}}
\newcommand{\dw}{\text{down}}

\newcommand{\diw}[1]{\widehat{#1}}
\newcommand{\di}[1]{\diw{#1}}
\newcommand{\bo}{b}
\newcommand{\ub}{u}
\newcommand{\fs}{\infty}
\newcommand{\ifst}{\infty}
\newcommand{\mon}{{mon}}
\newcommand{\cl}{\text{cl}}
\newcommand{\intr}{\text{int}}
\newcommand{\ul}[1]{\underline{#1}}

\newcommand{\bipil}{\leftrightarrow}
\newcommand{\bfc}{{\mathbf c}}

\renewcommand{\mp}{m^\prime}

\newcommand{\np}{n^\prime}
\newcommand{\Mod}{\text{Mod }}
\newcommand{\Sh}{\text{Sh } }
\newcommand{\st}{\text{st}}

\newcommand{\hM}{\tilde{M}}
\newcommand{\hs}{\tilde{s}}
\newcommand{\ee}{\mathbf{e}}

\newcommand{\lex}{{\text{lex}}}
\newcommand{\ordGL}{\succeq_{\lex}}
\newcommand{\ordG}{\succ_{\lex}}
\newcommand{\ordML}{\preceq_{\lex}}
\newcommand{\ordM}{\prec_{\lex}}
\newcommand{\tLa}{\tilde{\Lambda}}
\newcommand{\tGa}{\tilde{\Gamma}}
\newcommand{\STS}{\text{STS}}
\newcommand{\ii}{{\rotatebox[origin=c]{180}{\scalebox{0.7}{\rm{!}}}}}

\newcommand{\jj}{\mathbf{j}}

\newcommand{\shmod}{\texttt{shmod}\,}
\newcommand{\hf}{\underline{f}}
\newcommand{\Glim}{\lim}
\newcommand{\Gcolim}{\colim}
\newcommand{\fm}{f^{\underline{m}}}
\newcommand{\fn}{f^{\underline{n}}}
\newcommand{\gn}{g_{\mathbf{|}n}}
\newcommand{\se}[1]{\overline{#1}}
\newcommand{\cB}{{\mathcal{B}}}
\newcommand{\fin}{{\text{fin}}}
\newcommand{\Poset}{{\text{\bf Poset}}}
\newcommand{\Set}{{\text{\bf Set}}}

\newcommand{\Homi}{\text{Hom}^{L}}
\newcommand{\Homb}{\text{Hom}_S}
\newcommand{\Homu}{\text{Hom}^u}
\newcommand{\bfnu}{{\mathbf 0}}

\newcommand{\semip}{up semi-finite }
\newcommand{\semim}{down semi-finite }
\newcommand{\dif}[1]{\di{#1}_{fin}}
\newcommand{\llin}{\raisebox{1pt}{\scalebox{1}[0.6]{$\mid$}}}
\newcommand{\promap}{\mathrlap{{\hskip 2.8mm}{\llin}}{\lpil}}
\newcommand{\dual}{{\widehat{}}}
\newcommand{\Bool}{{\bf Bool }}
\newcommand{\pro}{{pro}}
\newcommand{\ovF}{\overline{F}}
\newcommand{\TO}{{\mathrm to}}
\newcommand{\dT}{\overset{\rightarrow}{T}}

\newlength{\dhatheight}
\newcommand{\doublehat}[1]{
    \settoheight{\dhatheight}{\ensuremath{\widehat{#1}}}
    \addtolength{\dhatheight}{-0.35ex}
    \widehat{\vphantom{\rule{1pt}{\dhatheight}}
    \smash{\widehat{#1}}}}

\newcommand{\napo}{natural }

\newcommand{\colim}{\text{colim}}
\newcommand{\oPsi}{\overline{\Psi}}
\newcommand{\bfC}{{\mathbf C}}
\newcommand{\bfi}{{\mathbf i}}
\newcommand{\ik}{{\mathfrak k}}
\newcommand{\ic}{{\mathfrak c}}

\newcommand{\ijfr}{{\mathfrak j}}
\newcommand{\iifr}{{\mathfrak i}}

\newcommand{\iC}{{\mathfrak C}}

\newcommand{\Ab}{A_\bullet}
\newcommand{\Rb}{R_\bullet}
\newcommand{\oDo}{\overline{D^\opos}}
\newcommand{\bfA}{{\mathbf A}}
\newcommand{\het}{\text{ht}}

\newcommand{\bfal}{{\mathbf{ \alpha}}}
\newcommand{\al}{{\alpha}}
\newcommand{\alp}{{\alpha^\prime}}
\newcommand{\tci}{tight}
\newcommand{\bitci}{bitight}
\newcommand{\oD}{\overline{D}}
\newcommand{\Delst}{\vardel \Delta}
\newcommand{\mlin}{-}

\newcommand{\ben}{{\mathbf 1}}
\newcommand{\bfj}{{\mathbf j}}
\newcommand{\bfio}{\hat{\bfi}}
\newcommand{\pot}{{\rm{pos}}}
\newcommand{\rac}{r-acyclic}

\newcommand{\Ign}[1]{}
\newcommand{\bfr}{{\mathbf r}}
\newcommand{\bfq}{{\mathbf q}}
\title{Artin monomial polarizations define full-dimensional simplicial balls}
\author{Gunnar Fløystad, Ine Gabrielsen, Amir Mafi}
\date{doi:10.5802/alco.500; CC BY 4.0}
\begin{document}
\maketitle
\noindent\textit{Editorial note.} Fixed official corrected published version only; no withdrawn bitight-equivalence import.\par\medskip
\noindent\textit{Editorial note.} Original author wording unchanged; counter-only formatting additions preserve numbers.\par\medskip
\noindent\textit{Editorial note.} Full published references reconstructed from official HTML with original numbers.\par\medskip
\noindent\textit{Editorial note.} No external image/diagram used in these selected spans; native author displays retained.\par\medskip
\noindent\textit{Editorial scope.} Faithful author material collection only. Root owns assembly, compile, visual QA, SQLite, upload and identity refresh.\par\medskip
\noindent\textit{Allowed standard prerequisites (editorial).}\begin{itemize}
\item Eagon–Reiner duality. Published [17]; quoted in Section2.4, native1215–1218. CM squarefree ideals correspond to Alexander dual ideals with linear resolution.
\item Björner constructible pseudomanifold criterion. Published [10], Theorem11.4; fully quoted as current Theorem4.2. Final conversion to ball or sphere after current constructibility proof.
\item CM quotients and regular sequences. Bruns–Herzog [13], Theorems2.1.3(a),2.1.2(c). Proposition3.6 proof retained, using these standard commutative-algebra tools.
\item Proper polarization extension bound. Almousa–Fløystad–Lohne [3], Remark2.5. External variable-count bound explicitly invoked by published proof; current structural characterization Proposition3.6 reproduced in full.
\item Multigraded resolutions and homological algebra. Standard Tor long exact sequence and mapping cone; published bibliography. Standard operations only; current truncation lemma and intersection/constructibility arguments are retained.
\end{itemize}\medskip
\section*{Source and counting boundary}
Gunnar Fløystad, Ine Gabrielsen and Amir Mafi, \emph{Polarizations of Artin monomial ideals}, Algebraic Combinatorics 9(4) (2026), 993--1035, \url{https://doi.org/10.5802/alco.500}.
This editable selection uses only the official published article and its publisher-provided native TeX, licensed CC BY 4.0: \url{https://creativecommons.org/licenses/by/4.0/}. Copyright remains with the authors.

Only the whole published Theorem 4.3 is counted. Its strict-containment hypothesis and proper-polarization extension clause are retained. Theorems 4.1--4.2, Proposition 3.6 and all supporting results are support, not extra items. Theorem 5.5 and the equivalent immediate consequence Theorem 4.4 are not counted.

\paragraph{Version boundary.} The earlier arXiv:2212.09528v2 was withdrawn on 5 August 2024: the authors reported an error in the proof of Lemma 6.17, 1a, Case iii, affecting old Theorem 6.7 (old part 3 does not imply parts 1 and 2). The v3 announcement of 13 January 2025 states that the error was fixed. This selection uses the actual 2026 published version alone, with its current linear-resolution/intersection/constructibility argument. No earlier bitight-equivalence text is imported and no mathematical repair or certification is supplied. Version notices: \url{https://arxiv.org/abs/2212.09528v2} and \url{https://arxiv.org/abs/2212.09528v3}.

\paragraph{Presentation.} Original mathematical wording, macros and selected equations are retained. Editorial counter settings reproduce published theorem and equation numbers across omitted unrelated blocks. The complete published reference list is rendered from the official landing-page bibliography with original numbers and matching native citation keys. The publisher class, logos and website styling are not reproduced. The retained proof contains no external image dependency.

\paragraph{Difficulty and responsibility.} The H2 core retained below is the multigraded truncation and Tor-degree control of all intersection ideals, followed by recursive constructibility over colour classes and the Alexander-dual topological conversion. This is a material collection, not an independent mathematical certification. The known in-flight identity snapshot ends on 14 October 2026 at 17:36:45 UTC; later overlap remains unconfirmed.

\setcounter{section}{1}\setcounter{cdrthm}{0}

\subsection{Separation and polarization}

For a field $k$ and a set $V$, let $k[V]$ be the polynomial ring with the
elements of $V$ as variables.
\begin{defi} \label{def:intro-sep}
Let $V^\prime \mto{p} V$ be a surjection of finite sets with the
cardinality of $V^\prime$ one more than that of $V$. Let $v_1$ and $v_2$
be the two distinct elements of $V^\prime$ which map to a single
element $v$ in $V$. Let $J$ be a monomial ideal in the polynomial ring
$k[V]$ and $J^\prime$ a monomial ideal in $k[V^\prime]$.
Then $J^\prime$ is a {\it simple separation} of $J$ if:
\begin{itemize}
\item[i.] The monomial ideal $J$ is the image of $J^\prime$ by the map
$k[V^\prime] \pil k[V]$.
\item[ii.] Both the variables $x_{v_1}$ and $x_{v_2}$ occur in some minimal
  generators of $J^\prime$ (usually in distinct generators, but in same
  generator is allowed).
\item[iii.] The variable difference $x_{v_1} - x_{v_2}$ is a non-zero divisor
in the quotient ring $k[V^\prime]/J^\prime$.
\end{itemize}

The minimal generators of $J$ and $J^\prime$ will then be
in bijection.
More generally, if $V^\prime \mto{p} V$ is a surjection of finite sets
and $J \sus k[V]$ and $J^\prime \sus k[V^\prime]$ are monomial ideals such
that $J^\prime$ is obtained by a succession of simple separations of $J$,
$J^\prime$ is a {\it proper separation} of $J$.
 If further $V^\prime \sus V^{\prime \prime}$ is an inclusion,
  the {\it extension} ideal $J^{\prime \prime} = (J^\prime) \sus k[V^{\prime \prime}]$
 is a {\it separation} of the starting ideal $J$.

\end{defi}

Note that $J^{\prime \prime}$ is obtained from $J$ by a sequence of
simple separations where we do not require ii. to hold. 

\begin{defi} \label{def:intro-pol}
  Let $J \sus k[V]$ be a monomial ideal and $V^\prime \pil V$ be a surjection
  of finite sets. An ideal $J^\prime \sus k[V^\prime]$ is a
  {\it polarization}
  of $J$ if $J^\prime$ is squarefree and a separation of $J$.
  It is a {\it proper polarization} if it is a squarefree proper
  separation of $J$. 
\end{defi}

\setcounter{equation}{0}

Given non-empty disjoint
finite sets $A_1, A_2, \ldots, A_n$ consider the simplicial join
\begin{equation*}
  \Delst(\Ab) = \pd \Delta(A_1) * \pd \Delta(A_2) * \cdots * \pd \Delta(A_n).
 \end{equation*}
 Its topological realization is the topological joins of the
 topological realizations of the $\pd \Delta(A_i)$,
 and so is a sphere of dimension $\sum_{i = 1}^n |A_i| - 1 - n$.
 In particular, if all $A_i$ have cardinality two, $\Delst(\Ab)$
 is the boundary of an $n$-dimensional hyperoctahedron.
 Let $A$ be the (disjoint) union of the $A_i$. We note that the facets
 of $\Delst(\Ab)$ are indexed by elements
 \begin{equation} \label{eq:intro-bfa}
   \bfa = (a_1, a_2, \ldots, a_n)  \in A_1 \times A_2 \times
   \cdots \times A_n.
   \end{equation}
 To such an element corresponds the facet $A \backslash \{a_1, \ldots, a_n\}$. 

  We relate these geometric objects to Artin monomial ideals.
 Let $k[A]$ be the
 polynomial ring with variables from $A$. Denote by $m(A_i)$
 the monomial $\prod_{a \in A_i} a$, and by $\alpha_i$  the cardinality of $A_i$. 
In the polynomial ring $k[x_1, \ldots, x_n]$ we have the complete
 intersection Artin monomial ideal:
 \begin{align*}
   \ic = (x_1^{\alpha_1}, x_2^{\alpha_2}, \ldots, x_n^{\alpha_n}).
 \end{align*}

\paragraph{Editorial context for the notation.} In the next source paragraph, such sub-complexes means full-dimensional Cohen--Macaulay proper subcomplexes of $\Delst(\Ab)$, as specified immediately before it in published Theorem 1.6.

In fact, such sub-complexes $X$ are precisely those that via the
Stanley-Reisner correspondence come from polarizations of Artin
monomial ideals. 
As noted above
  (see~\eqref{eq:intro-bfa}), the facets of
  $X$ are then indexed by a subset $T \sus A_1 \times \cdots \times A_n$.
  Thus we may write $X = X(T)$, and its Stanley-Reisner ideal $J = J(T)$.

  A standard fact (see Subsection~\ref{subsec:pre-const})
  is that $J$ is CM iff its Alexander dual $I$ is an ideal with linear
  resolution. We write $I(T)$ for the Alexander dual of $J(T)$. 
    Considering each $A_i$ as a color class, each
  $\bfa \in A_1 \times \cdots \times A_n$ gives
  a {\it rainbow monomial} $m(\bfa) = \prod_{i=1}^n a_i$.
  The ideal $I(T)$ is generated by the $m(\bfa)$ for $\bfa \in T$.

\setcounter{section}{1}\setcounter{subsection}{0}

\section{Preliminaries}
\label{sec:pre}
We recall basics on simplicial complexes,
the Stanley-Reisner correspondence, and Alexander duality.
For introduction to these, see~\cite{HeHi, MiSt} or~\cite[Sec.II.5]{BH93}.

\subsection{Simplicial complexes and squarefree monomial ideals}
Let $V$ be a finite set. A simplicial complex on $V$ is
a family $X$ of subsets of $V$ such that if $F \in X$
and $G \sus F$, then also $G \in X$. The $F$ are called the
{\it faces} of $X$ and the dimension of $F$ is the cardinality $|F| -1$.
Maximal faces are called {\it facets}. If all facets have the same
dimension, $X$ is said to be {\it pure}.
If $X$ is the full power set of $V$, then
$X$ is called the {\it simplex} on $V$ and is denoted $\Delta(V)$. 
The {\it boundary $\pd X$} of a pure
simplicial complex $X$ is the smallest simplicial
complex containing all faces of $X$ which are on only one facet, and
are properly contained in this facet.
A simplicial complex $X$ has a natural topological realization $|X|$ in the
space ${\mathbb R}^V$,~\cite{Mun}.

Let $k$ be a field and $k[V]$ the polynomial ring with elements of $V$
as variables. If $R \sus V$ denote by $m(R) = \prod_{v \in R} v$ the
associated monomial. Such a monomial is said to be {\it squarefree},
and if an ideal is generated by such monomials, it is also said to be
{\it squarefree}.



Given a simplicial complex $X$ we have its associated Stanley-Reisner (SR)
ideal $I= I_X$ which is the ideal generated by the monomials
$m(R)$ such that $R$ is {\it not} in $X$. The quotient ring $k[V]/I$ is
the Stanley-Reisner ring of $X$.

Simplicial complexes form a lattice with join $X \cup Y$, the union
of the families $X$ and $Y$, and meet $X \cap Y$.
The squarefree monomial ideals also form a lattice
with join the sum of ideals, and meet the intersection of ideals.

The Stanley-Reisner correspondence gives
a one-to-one correspondence between the lattice of
simplicial complexes on $V$
and the lattice of squarefree monomial ideals in $k[V]$, with order relation
reversed. Hence the join $X \cup Y$ corresponds to the meet of ideals
$I_X \cap I_Y$, and the meet $X \cap Y$ corresponds to the sum of ideals
$I_X + I_Y$. 


\subsection{Alexander duality}
\label{subsec:AD}
The Alexander dual monomial ideal of $I (=I_X)$
is the monomial ideal $J = I^\vee$ generated
by all monomials $m$ which have a non-trivial common divisor with every monomial
in $I$. This is an involutive correspondence. The
squarefree monomials in $J$ are the $m(F^c)$ where the $F$ are faces
of $X$, and $F^c = V \backslash F$ its complement. In particular the
minimal generators
of $J$ are the monomials $m(F^c)$ of the complements of facets $F$
of $X$, see~\cite[Lemma 1.5.3]{HeHi}.

Alexander duality gives a bijection between lattices of ideals
with order relation reversed. Hence the Alexander dual of the join
$I + L$ is the meet $I^\vee \cap L^\vee$, and the dual of $I \cap L$ is
$I^\vee + L^\vee$.


\subsection{External join of simplicial complexes}

Let $X$ be a simplicial complex on~$V$, and $Y$ a simplicial complex
on a disjoint vertex set $W$. The (external) join of~$X$ and $Y$ is
the simplicial complex $X * Y$ on $V \cup W$ given by:
\[ X * Y = \{ A \cup B \, | \, A \in X, B \in Y \}. \]
(This is not to be confused with the (internal) join in the lattice
of simplicial
complexes on a fixed set $V$, which is simply the union of $X$ and $Y$.
Henceforth join will always mean external join.)
The SR-ideal of the join $X * Y$ is the ideal
$(I_X) + (I_Y)$ generated by $I_X$ and $I_Y$ in $k[V \cup W]$. 

The topological realization $|X * Y|$ of the join, is obtained by connecting
points
in $|X|$ and $|Y|$ by lines. So it is the set
$|X| \times [0,1] \times |Y|$ modulo identifying pairs
$(x,0,y)$ and $(x,0,y^\prime)$ and identifying pairs
$(x,1,y)$ and $(x^\prime, 1, y)$.
The join of two spheres $S^n$ and $S^{m}$ of dimensions $n$ and $m$,
is the sphere $S^{n+m+1}$. If $|Y|$ is two points (corresponding
to the simplicial complex $\{\{v_1\}, \{v_2\}, \emptyset\}$, which
is topologically the sphere $S^0$), the
join $X * Y$ has topological realization the suspension $S|X|$.
The suspension of $S^n$ is $S^{n+1}$.

\subsection{Constructible squarefree monomial ideals}
\label{subsec:pre-const}
A simplicial complex of dimension $d$ is {\it constructible} if either
\begin{itemize}
\item $X$ is a simplex of dimension $d$, or
\item there is a decomposition $X = Y \cup Z$ where $Y$ and $Z$
  are constructible of dimension $d$ and $Y \cap Z$ is constructible
  of dimension $d-1$.
\end{itemize}
The following implications are well known:
\[ \text{vertex decomposable} \Rightarrow
  \text{shellable} \Rightarrow \text{constructible} \Rightarrow
  \text{Cohen-Macaulay/any field}. \]
Let $I$ be the Alexander dual ideal of the Stanley-Reisner ideal $I_X$.
Translating the above to conditions on the ideal $I$ we get the following.
A squarefree monomial ideal $I$ is {\it $e$-constructible}
(notion introduced in~\cite{Olt-Const}) if:
\begin{itemize}
\item $I = (m)$ is generated by a monomial of degree $e$, or
\item there is a decomposition $I = J + K$ such that $J$ and $K$
  are $e$-constructible, and $J \cap K$ is $(e+1)$-constructible.
\end{itemize}

\begin{rema} If $I$ corresponds to the simplicial complex $Y$ via
  the SR-correspondence, one should note that $I$ being constructible
  does {\it not} correspond to $Y$ being constructible. Rather $I$ being
  constructible corresponds to 
  the Alexander dual simplicial complex $Y^\vee$ being constuctible
  (where  $Y^\vee$ has the Alexander dual ideal $I^\vee$ as
  Stanley-Reisner ideal).
\end{rema}

The Eagon-Reiner theorem,~\cite{ER98}, states that a squarefree
monomial ideal $I$ has a linear resolution if and only if the Alexander dual $I^\vee$
is a Cohen-Macaulay ideal, see standard textbooks~\cite[Thm.~8.1.9]{HeHi},
\cite[Thm.~5.56]{MiSt}.
As a consequence constructible ideals have linear resolutions.
In general for ideals there are implications:
\[  \text{vertex splittable} \Rightarrow
  \text{linear quotients} \Rightarrow \text{constructible} \Rightarrow
  \text{linear resolution/any field}. \]
The notion of vertex splittable is introduced in~\cite{MoKa}, and
linear quotients in~\cite{HT-cones}.

\setcounter{section}{3}\setcounter{subsection}{0}\setcounter{cdrthm}{0}\setcounter{equation}{3}

\subsection{Joins of boundaries of simplices}
Recall that $\Delta(V)$ is the simplex on the finite set $V$. 
The SR-ideal generated by the single monomial $m(V) = \prod_{v \in V} v$
defines the
boundary $\pd \Delta(V)$ of $\Delta(V)$, which is topologically
a sphere $S^{|V|-2}$ where $|V|$ is the cardinality of $V$. 

Let  $A_1, A_2, \ldots, A_n$ be disjoint finite sets, and $A = \sqcup_{i=1}^n A_i$
their union. The complete intersection ideal
\begin{equation*}
  C = (m(A_1), m(A_2), \ldots, m(A_n))  \sus k[A]
\end{equation*}
is the SR-ideal of the join
of the boundaries of simplices
\begin{equation} \label{eq:liaison-join}
  \Delst(\Ab) = \pd \Delta(A_1) * \pd \Delta(A_2) * \cdots * \pd \Delta(A_n),
 \end{equation}
which topologically is a join of spheres, and so also a sphere.
Its dimension is $\sum_{i = 1}^n |A_i| -1-n$. If each cardinality
$|A_i| = 2$, then each boundary $\pd \Delta(A_i)$ is two points,
and we are taking iterated suspensions. This gives a sphere
of dimension $n-1$: For $n = 2$, (the boundary of) the square,
for $n = 3$
(the boundary of) the octahedron, and in general
(the boundary of) the hyperoctahedron.

\subsection{Rainbow ideals}
Denote $\bfA = A_1 \times A_2 \times \cdots \times A_n$.
For $\bfa = (a_1, a_2, \ldots, a_n)$ in $\bfA$, denote
by $[\bfa ]$ the set $\{a_1, a_2, \ldots, a_n \}$. 
The faces of $\Delst(\Ab)$ are all subsets $S \sus A$ such
that no $A_i \sus S$. In other words it is all $S$ such that
each $A_i \backslash S$ is non-empty. Equivalently we 
may find $\bfa \in \bfA$ with $[\bfa] \cap S = \emptyset$.
The facets of $\Delst(\Ab)$ are then precisely
the complement sets $[\bfa]^c$ for $\bfa \in \bfA$.

The following is of course true for any triangulated sphere, but the argument
is explicit and informative in our setting.

\begin{lemm} \label{lem:link-cdto}
  Any codimension one face of $\Delst(\Ab)$ is on exactly two facets.
\end{lemm}

\begin{proof} A codimension one face of $\Delst(\Ab)$ is
  $F = [\bfa]^c \backslash \{a^\prime \}$ for some $a^\prime \in A$.
  So there is an $i$ such that $a^\prime \in A_i$, and write $a_i^\prime
:= a^\prime$. Let $\bfa^\prime = (a_1, \ldots, a_i^\prime, \ldots, a_n)$.
  Then $F$ is precisely on the two facets $[\bfa]^c$ and $[\bfa^\prime]^c$.
\end{proof}


For $\bfa \in \bfA$ let $m(\bfa)$ be the monomial 
$\prod_{i = 1}^n a_i$. Thinking
of $A_i$ as a color class, we call such a monomial a {\it rainbow
monomial}.  The Alexander dual ideal $B$ of the complete intersection
$C = (m(A_1), \ldots, m(A_n))$ is generated by the complements
of the facets of $\Delst(\Ab)$, and so it is the monomial
ideal generated by all the  rainbow monomials. Letting
$(A_i)$ be the ``color ideal'' generated by the elements of $A_i$ (considered
as variables), then $B$ is the product of these color ideals:
\[ B  = (m(\bfa) \, | \, \bfa \in \bfA) = \prod_{i = 1}^n(A_i). \]

\begin{defi}
For a subset $T \sus \bfA$, denote by $I(T)$ the ideal generated
by $m(\bfa)$ for $\bfa \in T$. Such an ideal is called a {\it rainbow ideal}.
For the complement set $T^c = \bfA \backslash T$, we
usually write $L(T)$ for $I(T^c)$, the {\it complement} rainbow ideal.
\end{defi}

\setcounter{subsection}{2}\setcounter{cdrthm}{5}

  \subsection{Polarizations}
  \label{subsec:pol}

  Recall from the introduction, Definitions~\ref{def:intro-sep} and
 ~\ref{def:intro-pol} of separations and polarizations of monomial ideals.

  Let $\ijfr$ be an Artin monomial ideal
  in $k[x_1, \ldots, x_n]$, containing the Artin complete
  intersection
  \[ \ic = (x_1^{\al_1}, \ldots, x_n^{\al_n}). \]
  We do not require the powers $x_i^{\alpha_i}$ 
  to be minimal generators of the ideal $\ijfr$. 
  Let $A_1, \ldots, A_n$ be disjoint sets with $A_i$ of cardinality $\alpha_i$,
  and $A = \cup_i A_i$.
We have a map of polynomial rings
  $k[A] \pil k[x_1, \ldots, x_n]$ induced by the map 
  $A \pil \{x_1,x_2, \ldots, x_n \}$ sending elements of $A_i$
  to the variable $x_i$.
  By~\cite[Remark 2.5]{AFL} any proper polarization $J^\prime$
  of an Artin monomial ideal
  containing the complete intersection $\ic$ is, after possibly
  extending $J^\prime$, isomorphic to a polarization $J$ contained in $k[A]$. 
  
  We also recall the following, which is essentially~\cite[Prop.3.6]{AFL}.
  Due to its importance here,
  we include the proof for the benefit of the reader.
  
\begin{prop} \label{pro:pol-cm}
  Let $C = (m(A_1), \ldots, m(A_n))$ be the complete intersection.
  Polarizations $J \sus k[A]$ of Artin monomial ideals containing the complete
  intersection $\ic$
  are precisely the Cohen-Macaulay
  squarefree ideals $J \sus k[A]$ of height $n$  containing  $C$.

  Such ideals $J$ are in turn precisely 
    the Alexander duals of rainbow monomial ideals in $k[A]$ with
    linear resolution.
    \end{prop}

    \begin{proof} 
      i. Assume $J \subseteq k[A]$ is a polarization of an Artin
      monomial ideal containing~$\ic$. Then $J$ contains $C$, as
      $m(A_i)$ must be the monomial mapping to $x_i^{\al_i}$.
      Since $C$ is cut down to an Artin monomial ideal, $J$ is also cut
      down by the regular sequence (for $k[A]/J$) to an Artin monomial ideal.
      An Artin ideal is Cohen-Macaulay, and so the same holds for $J$ by
     ~\cite[Thm.~2.1.3.a]{BH93}.

      ii. Conversely assume $J$ is Cohen-Macaulay and squarefree of height $n$
      containing~$C$. The sequence
      of variable differences cuts $C$ down to an Artin monomial ideal,
      and so also cuts $J$ down to an Artin monomial ideal $\ijfr$. Thus
      starting from $k[A]/J$, which has dimension $|A|-n$, 
      at each step when we divide out by a variable difference,
      dimension is cut by one. By~\cite[Thm.~2.1.2.c]{BH93}
      the sequence is regular for $k[A]/J$ 
      and so $J$ is a polarization of $\ijfr$. 
      
      By the Eagon-Reiner theorem, see Subsection
     ~\ref{subsec:pre-const},
      $J$ is Cohen-Macaulay iff its Alexander dual $I$ has a linear resolution.
      Let $X$ be the simplicial complex associated to $J$. That $J$ is
      unmixed of height $n$ is equivalent to  the facets of $X$ being
      facets of $\Delst(\Ab)$, and so $I = I(T)$ for some $T \sus \bfA$.
    \end{proof}

\setcounter{section}{3}\setcounter{subsection}{0}

\section{Polarizations define triangulated balls}
\label{sec:polball}

The previous section shows that
polarizations of
Artin monomial ideals in $k[x_1, \ldots, x_n]$
containing the complete intersection $\ic = (x_1^{\alpha_1},
\ldots, x_n^{\alpha_n})$
correspond, by Alexander duality, to rainbow ideals in $k[A]$ with linear resolutions,
for $A$ a union of color classes
$A_1, \ldots, A_n$ and cardinalities $|A_i|= \alpha_i$.
Rainbow ideals
are $I(T)$ for a subset $T \sus \bfA = A_1 \times \cdots \times A_n$.

We investigate ideals of type
\[ I = b_1 I_1 + b_2 I_2 + \cdots + b_m I_m \]
where the $I_j$ do not involve the $b$-variables, and the situation
when such ideals have a linear resolution.
We apply this to show that rainbow ideals have
linear resolutions if and only if they
are constructible. As a consequence we get our main theorem,
that polarizations of Artin monomial ideals define triangulated
balls by the Stanley-Reisner correspondence.

\subsection{The main result}

The following is the key result leading to our main statement.
  
 \begin{theo} \label{thm:linres-const} For a subset
   $T \sus A_1 \times \cdots \times A_n$,
   the rainbow ideal $I(T)$ has $n$-linear resolution if and
   only if it is $n$-constructible.
  \end{theo}

  By Subsection~\ref{subsec:pre-const}, if  $I(T)$ is constructible, then
  $I(T)$ has a linear resolution. The challenge is to show the other direction.
  But let us first give the consequence which is our main result.
  First recall the following, stated in the overview article
 ~\cite[Thm.~11.4]{Bj95} by A.~Bj\"orner:


\begin{theo} \label{thm:bjorner}
Let $X$ be a pure constructible simplicial complex
    such that every codimension one face is on one or two facets.
    Then the topological realization $|X|$, is a ball or a sphere
    (the latter when every codimension one face is on exactly two
    facets).
\end{theo}
  
Using this, Theorem~\ref{thm:linres-const} implies Conjecture 3.2 in~\cite{AFL}.


\begin{theo}[Main Theorem] \label{thm:rainbow-main}
Let $\ijfr$ be an Artin monomial ideal in $k[x_1, \ldots, x_n]$
  strictly containing the
  complete intersection $\ic$ (but the $x_i^{\alpha_i}$ need not be
  minimal generators of $\ijfr$).
  Let $A \pil \{x_1, x_2, \ldots , x_n\}$  have
  fiber $A_i$ over $x_i$ of cardinality $\alpha_i$.
  
  Any polarization 
  $J \sus k[A]$ of $\ijfr$ defines, via the SR-correspondence,
  a  full dimensional  simplicial ball on the simplicial sphere $\Delst(\Ab)$.
  Moreover, any proper polarization of $\ijfr$ is,
  when followed by
  a suitable extension,
  isomorphic to such a $J$. 
\end{theo}

\begin{proof}
  By~\cite[Remark 2.5]{AFL}, any proper polarization of $\ijfr$ will have
  at most $\alpha_i$ variables mapping to $x_i$. Hence it is, after
  possible extension, isomorphic to a polarization $J \sus k[A]$. 
  By Proposition~\ref{pro:pol-cm} such a polarization is of the
  form $J = J(T)$, and is Cohen-Macaulay. Denote the
  associated  simplicial complex by $X(T)$. 
  The Alexander dual  $I(T)$ of $J(T)$ has a linear resolution. By Theorem
 ~\ref{thm:linres-const}, $I(T)$ is constructible and so $X(T)$ is
  constructible. By Lemma~\ref{lem:link-cdto} every codimension one
  face of $X(T)$ is on at most two facets. Theorem~\ref{thm:bjorner} above
  gives that $X(T)$ is topologically a ball.
\end{proof}

\setcounter{cdrthm}{4}\setcounter{equation}{4}

\subsection{Ideals with linear resolutions and constructability}
 
We proceed to prove Theorem~\ref{thm:linres-const}.
First we give a general fact. Let $k[B \cup X]$ be a polynomial ring,
and $\hele^B$ the free abelian group with basis $e_b$ for $b \in B$.
We consider the polynomial ring to be $\hele^B \oplus \hele$ graded,
where the variables  $b \in B$ have degree $e_b$ and the variables
in $X$ have degree $e = (\mathbf 0,1)$.


The following is a crucial fact: If $p$ is a homogeneous
polynomial in $k[B \cup X]$ (for the grading by $\ZZ^B \oplus \ZZ$),
of degree $(\bfr,s)$, then $p$ is a product
$(\prod_{b \in B}b^{r_b}) \cdot p^\prime$ where $p^\prime$ is a homogeneous
polynomial in $k[X]$ of degree $s$.

Let $M$ be a $\ZZ^B \oplus \ZZ$-graded 
$k[B \cup X]$-module. If $m \in M$ is homogeneous of degree $(\bfr,s)$,
write $\deg_B(m) = \bfr$.
For $\bfr \in \ZZ^B$, denote $M_{\leq \bfr}$ be the
submodule of $M$ generated by the homogeneous elements $m \in M$
with $\deg_B(m) \leq \bfr$.
Note that if $F = \oplus_{j \in J} S(-{\mathbf u}_j)$ is a free module,
then $F_{\leq \bfr} = \oplus_{\deg_B({\mathbf u}_j) \leq \bfr} S(-{\mathbf u}_j)$.

\begin{lemma} Let $\phi : N \pil M$ be a morphism of $\ZZ^B \oplus \ZZ$-graded
  $k[B \cup X]$-modules, with $M$ a torsion free module. Let $\bfr \in \ZZ^B$. 
  If $K$ is the kernel of $\phi$, then $K_{\leq \bfr}$ is the kernel of
  $\phi_{\leq r} : N_{\leq \bfr} \pil M_{\leq \bfr}$.
\end{lemma}

\begin{proof} Clearly $K_{\leq \bfr}$ is contained in the kernel of
  $\phi_{\leq \bfr}$. 
  Suppose $n$ is homogeneous and
  in the kernel of $\phi_{\leq \bfr}$. We may write 
  $n = \sum p_i n_i$, with $p_i$,$n_i$ homogeneous in respectively $k[B \cup X]$ and $N$,
  and with $\deg_B(n_i) \leq \bfr$. Let $\deg_B(n) = \bfq$, 
  and consider the monomial
  $u = \prod_{q_b > r_b}b^{q_b - r_b}$. Since each $\deg_B(n_i) \leq \bfr$,
  we have each $\deg_B(p_i) \geq \deg(u)$. But then $p_i$ is
  divisible by~$u$. Thus we may write $n = u \cdot n^\prime$
  with $\deg_B(n^\prime) \leq \bfr$. 

  Then $0 = \phi(n) = u \phi(n^\prime)$ and since $M$ is torsion
  free, this gives $\phi(n^\prime) = 0$, and so $n^\prime \in \ker \phi = K$.
  Thus  $n \in K_{\leq \bfr}$.
\end{proof}

This gives the following (a generalization of~\cite[Thm.~2.1]{GaHiPe}).

\begin{coro}
  Let $M$ be a $\ZZ^B \oplus \ZZ$-graded torsion free $k[B \cup X]$-module,
  and $F_\bullet$ a minimal free resolution of $M$. Let $\bfr \in \ZZ^B$.
  Then $F_{\bullet, \leq \bfr}$ is a minimal free resolution of~$M_{\leq \bfr}$.
\end{coro}

We say $F_{\bullet}$ is a {\it $d$-linear resolution}, if
$F_i = \oplus_{j \in J} S(-{\mathbf u}_{ij})$ and for each
$\deg(u_{ij}) = \bfq \oplus s$, its total degee $s + \sum_{b \in B} q_b$
equals $d + i$.

\begin{coro} \label{cor:lin-Frlin}
  If $F_{\bullet}$ is a $d$-linear resolution of $M$, then
  $F_{\bullet, \leq \bfr}$ is a $d$-linear resolution of $M_{\leq \bfr}$.
\end{coro}

In the following we let $B = \{b_1, b_2, \ldots, b_m\}$ have
cardinality $m$. We write $\ZZ^B = \ZZ^m$ and $e_i = e_{b_i}$ for
$i = 1, \ldots, m$, and $e_{m+1} = {\mathbf 0} \oplus 1$.

\begin{prop} \label{pro:lin-ItoIR}
  Let $I_1, \ldots, I_m$ be ideals in $k[X]$.
Consider the product ideals $(b_j) (I_j) \sus k[B \cup X]$, which
we simply denote as $b_j I_j$. Also for $R \sus \{1,2, \cdots, m\}$
denote
\[ I_R = \cap_{j \in R} I_j  \sus k[X]. \]
If the ideal
\[ I = b_1 I_1 + b_2 I_2 + \cdots + b_m I_m \]
has $(d+1)$-linear resolution, then each $I_R$ has $d$-linear resolution.
\end{prop}

\begin{proof}
  First we claim: For each   $R \sus \{1,2, \ldots, m\}$
  the sub-sum  $\sum_{j \in R} b_j I_j$ has $(d+1)$-linear resolution.
  Let $\bfr = \sum_{j \in R} e_j $. Then $I_{\leq \bfr}$ is 
 $\sum_{j \in R} b_j I_j$, and by Corollary~\ref{cor:lin-Frlin}
it has a linear resolution. 
As a special case each  $b_j I_j$ has $(d+1)$-linear
resolution, and so each $I_j$ has $d$-linear resolution.


We now use induction on $m$. If $m = 1$ the statement holds.
Let $m = 2$. We have
an exact sequence
\[ 0 \pil b_1 I_1 \cap b_2I_2 \pil b_1I_1 \oplus b_2I_2 \pil
b_1 I_1 + b_2 I_2 \pil 0, \]
and a long exact $\Tor$-sequence:
{\small \begin{align*}  \cdots \pil & \Tor_{p+1}(b_1I_1 + b_2 I_2, k)_e \pil
  \Tor_{p}(b_1b_2 (I_1 \cap I_2), k)_e \pil \Tor_p(b_1I_1,k)_e \oplus
                           \Tor_p(b_2I_2,k)_e \\
  \pil & \Tor_{p}(b_1I_1 + b_2 I_2, k)_e \pil \cdots .
\end{align*}}Since $I = b_1I_1 + b_2 I_2$ has $(d+1)$-linear resolution, the Tor-groups
to the left vanish for $e - (p+1) \geq d+2$ or equivalently for $e-p \geq d+3$.
Similarly, to the right above the Tor-groups vanish for $e-p \geq d+2$.
Thus
\[ \Tor_{p}(b_1b_2 (I_1 \cap I_2), k)_e = 0, \quad e-p \geq d+3.\]
As $I_1 \cap I_2$ is generated in degree $\geq d$ it must have a
$d$-linear resolution.

Assume now $m = 3$. From
\[ 0 \pil (b_1 I_1 +  b_2I_2) \cap b_3I_3 \pil (b_1I_1 + b_2I_2)
  \oplus b_3 I_3 \pil
b_1 I_1 + b_2 I_2 + b_3I_3 \pil 0, \]
the same argument gives that $(b_1I_1 + b_2 I_2) \cap (I_3)$ has
$(d+1)$-linear resolution. Let $b_1 i_1 + b_2 i_2$ be a generator for
this ideal in degree $d+1$. It is also in $(I_3)$, but since $I_3$ is generated
in degree $d$ and its generators do not involve the variables $b_1$ and $b_2$,
both $i_1, i_2 $ are in $(I_3)$. 
Therefore
\[ (b_1 I_1 + b_2I_2) \cap (I_3) = b_1 I_1 \cap I_3 + b_2 I_2 \cap I_3. \]
By the $m = 2$ case, both $I_1 \cap I_3$ and $I_2 \cap I_3$ have
$d$-linear resolution, as well as their intersection $I_1 \cap I_2 \cap I_3$.

For larger $m$ we split as $(b_1I_1 + \cdots + b_{m-1}I_{m-1}) + b_mI_m$, and
proceed in the same way as above.
\end{proof}

The converse of Proposition~\ref{pro:lin-ItoIR} is not true
in general, but for monomial ideals it is.
We now proceed to do this. The key is the following:

\begin{lemma}
  Let $I$ and $J_1, \ldots, J_m$ be monomial ideals.
  Then
  \[ I \cap (J_1 + \cdots + J_m) = I \cap J_1 + \cdots + I \cap J_m. \]
\end{lemma}

\begin{proof}
  It is enough to prove this for $m = 2$, as the general case easily follows
  by induction. Suppose $f \in I \cap (J_1 + J_2)$. Then every monomial
  in $f$ is in $I$ and in $J_1 + J_2$. Then the monomial is in either
  $J_1$ or $J_2$. Whence $f = f_1 + f_2$ where $f_1$ consists of monomials
  in $I \cap J_1$ and $f_2$ of monomials in $I \cap J_2$.
\end{proof}

\begin{prop}
  Let $I_1, \ldots, I_m$ be monomial ideals in $k[X]$ such that for every
  $R \sus [m]$, the ideal $I_R = \cap_{j \in R} I_j$ has $d$-linear resolution.
  Then the ideal $I = \sum_{j=1}^n b_jI_j$ in $k[B \cup X]$ has
  $(d+1)$-linear resolution.
\end{prop}

\begin{proof}
  Again we use induction on $m$. Write
  \begin{equation} \label{eq:linres-KI}
    I = (b_1I_1 + \cdots + b_{m-1}I_{m-1}) + b_m I_m = K_{m-1} + b_mI_m.
    \end{equation}
    By induction we may assume $K_{m-1}$ has $(d+1)$-linear resolution. Note
    also that 
  $I_m$ has $d$-linear resolution. We have a short exact sequence:
  \begin{equation} \label{eq:linres-KIsekv}
    0 \pil K_{m-1} \cap b_mI_m \overset{f}\pil K_{m-1}
      \oplus b_m I_m \pil I \pil 0.
\end{equation}      
For the first term:
\begin{align*}
  K_{m-1} \cap b_mI_m&  = (b_1I_1 + \cdots + b_{m-1}I_{m-1}) \cap b_m I_m \\
                     &  = b_m(b_1I_1 \cap I_m + \cdots + b_{m-1} I_{m-1} \cap I_m)
\end{align*}
By induction the last ideal above has $(d+2)$-linear resolution. So
$K_{m-1} \cap b_m I_m$ has $(d+2)$-linear resolution. The mapping cone
of $f$ in the exact sequence~\eqref{eq:linres-KIsekv}
then becomes a $(d+1)$-linear resolution of $I$.
\end{proof}

In fact, almost the same argument as above may be applied for
the notion of constructability.

\begin{prop} \label{pro:lin-constr}
  Let $I_1, \ldots, I_m$ be monomial ideals in $k[X]$. If for every
  $R \sus [m]$, $I_R$ is $d$-constructible, then
  \[ I = b_1 I_1 + \cdots + b_m I_m \]
  is $(d+1)$-constructible.
\end{prop}

\begin{proof}
  We write $I = K_{m-1} + b_mI_m$ as in~\eqref{eq:linres-KI}.
  Here $b_m I_m$ is $(d+1)$-constructible
  and by induction on $n$, $K_{m-1}$ is also $(d+1)$-constructible.
  Further their intersection is:
  \begin{align*}  K_{m-1} \cap b_mI_m
    & = (b_1I_1 + \cdots + b_{m-1}I_{m-1}) \cap b_m I_m \\
    & = b_m(b_1I_1 \cap I_m + \cdots + b_{m-1} I_{m-1} \cap I_m),
\end{align*}
  and by induction, this is $(d+2)$-constructible.
  Thus $I$ is $(d+1)$-constructible.
\end{proof}

\begin{proof}[Proof of Theorem~\ref{thm:linres-const}.]
  When we have one color class,  $n = 1$, $I$ is generated by variables and
  so is constructible. Let $n \geq 2$, and $A_1 = \{a_1, a_2, \ldots, a_m \}$.
   We may then write
  \[ I = a_1I_1 + \cdots + a_m I_m \] where the $I_j = I(T_j)$ for
  various $T_j \sus A_2 \times \cdots \times A_n$. Since $I$ has $n$-linear
  resolution, for every $R \sus [m]$, $I_R = \cap_{j \in R} I_j$ has
  $(n-1)$-linear resolution. As all the $I_j$ are monomial ideals
  generated in degree $(n-1)$, we must have $I_R = I(T_R)$ where
  $T_R = \cap_{j \in R} T_j$. By induction on $n$, each of these are
  $(n-1)$-constructible. Then by Proposition~\ref{pro:lin-constr}, $I$ is
  $n$-constructible.
\end{proof}

% Full official published reference list; original numeric labels.
\begin{thebibliography}{51}
\bibitem[1]{published1} Aguiar, Marcelo; Ardila, Federico Hopf monoids and generalized permutahedra, Mem. Amer. Math. Soc., Volume 289 (2023) no. 1437, p. vi+119. \url{https://doi.org/10.1090/memo/1437}.
\bibitem[2]{published2} Alilooee, Ali; Banerjee, Arindam Symbolic powers of multipartite hypergraphs and Waldschmidt constant, 2021. \url{https://arxiv.org/abs/2103.06468}.
\bibitem[3]{AFL} Almousa, Ayah; Fløystad, Gunnar; Lohne, Henning Polarizations of powers of graded maximal ideals, J. Pure Appl. Algebra, Volume 226 (2022) no. 5, Paper no. 106924, 33 pages. \url{https://doi.org/10.1016/j.jpaa.2021.106924}.
\bibitem[4]{published4} Almousa, Ayah; VandeBogert, Keller Polarizations and hook partitions, J. Pure Appl. Algebra, Volume 226 (2022) no. 9, Paper no. 107056, 20 pages. \url{https://doi.org/10.1016/j.jpaa.2022.107056}.
\bibitem[5]{published5} Beckenbach, Isabel; Scheidweiler, Robert Perfect $f$-matchings and $f$-factors in hypergraphs—a combinatorial approach, Discrete Math., Volume 340 (2017) no. 10, pp. 2499-2506. \url{https://doi.org/10.1016/j.disc.2017.05.005}.
\bibitem[6]{published6} Beckenbach, Isabel Leonie Matchings and flows in hypergraphs, Ph. D. Thesis, FU-Berlin (2019).
\bibitem[7]{published7} Berkesch, Christine; Erman, Daniel; Smith, Gregory G. Virtual resolutions for a product of projective spaces, Algebr. Geom., Volume 7 (2020) no. 4, pp. 460-481. \url{https://doi.org/10.14231/ag-2020-013}.
\bibitem[8]{published8} Berkesch, Christine; Klein, Patricia; Loper, Michael C.; Yang, Jay Homological and combinatorial aspects of virtually Cohen-Macaulay sheaves, Trans. London Math. Soc., Volume 8 (2021) no. 1, pp. 413-434. \url{https://doi.org/10.1112/tlm3.12036}.
\bibitem[9]{published9} Bier, Thomas A remark on Alexander duality and the disjunct join, 1992 (preprint, 8 pages).
\bibitem[10]{Bj95} Björner, A. Topological methods, Handbook of combinatorics, Vol. 1, 2, Elsevier Sci. B. V., Amsterdam, 1995, pp. 1819-1872.
\bibitem[11]{published11} Björner, Anders; Paffenholz, Andreas; Sjöstrand, Jonas; Ziegler, Günter M. Bier spheres and posets, Discrete Comput. Geom., Volume 34 (2005) no. 1, pp. 71-86. \url{https://doi.org/10.1007/s00454-004-1144-0}.
\bibitem[12]{published12} Bruns, Winfried; Gubeladze, Joseph Polytopes, rings, and ${K}$-theory, Springer Monographs in Mathematics, Springer, Dordrecht, 2009, xiv+461 pages. \url{https://doi.org/10.1007/b105283}.
\bibitem[13]{BH93} Bruns, Winfried; Herzog, Jürgen Cohen-Macaulay rings, Cambridge Studies in Advanced Mathematics, 39, Cambridge University Press, Cambridge, 1993, xii+403 pages.
\bibitem[14]{published14} D’Alì, Alessio; Fløystad, Gunnar; Nematbakhsh, Amin Resolutions of co-letterplace ideals and generalizations of Bier spheres, Trans. Amer. Math. Soc., Volume 371 (2019) no. 12, pp. 8733-8753. \url{https://doi.org/10.1090/tran/7560}.
\bibitem[15]{published15} De Loera, Jesús A.; Rambau, Jörg; Santos, Francisco Triangulations: Structures for algorithms and applications, Algorithms and Computation in Mathematics, 25, Springer-Verlag, Berlin, 2010, xiv+535 pages. \url{https://doi.org/10.1007/978-3-642-12971-1}.
\bibitem[16]{published16} Dufossé, Fanny; Kaya, Kamer; Panagiotas, Ioannis; Uçar, Bora Effective heuristics for matchings in hypergraphs, Analysis of experimental algorithms (Lecture Notes in Comput. Sci.), Volume 11544, Springer, Cham, 2019, pp. 248-264. \url{https://doi.org/10.1007/978-3-030-34029-2_17}.
\bibitem[17]{ER98} Eagon, John A.; Reiner, Victor Resolutions of Stanley-Reisner rings and Alexander duality, J. Pure Appl. Algebra, Volume 130 (1998) no. 3, pp. 265-275. \url{https://doi.org/10.1016/S0022-4049(97)00097-2}.
\bibitem[18]{published18} Eisenbud, David Commutative algebra: With a view toward algebraic geometry, Graduate Texts in Mathematics, 150, Springer-Verlag, New York, 1995, xvi+785 pages. \url{https://doi.org/10.1007/978-1-4612-5350-1}.
\bibitem[19]{published19} Ene, Viviana; Herzog, Jürgen; Mohammadi, Fatemeh Monomial ideals and toric rings of Hibi type arising from a finite poset, European J. Combin., Volume 32 (2011) no. 3, pp. 404-421. \url{https://doi.org/10.1016/j.ejc.2010.11.006}.
\bibitem[20]{published20} Fløystad, Gunnar Partitions of vertices and facets in trees and stacked simplicial complexes, Graphs Combin., Volume 40 (2024) no. 4, Paper no. 79, 22 pages. \url{https://doi.org/10.1007/s00373-024-02804-6}.
\bibitem[21]{published21} Fløystad, Gunnar; Møller Greve, Bjørn; Herzog, Jürgen Letterplace and co-letterplace ideals of posets, J. Pure Appl. Algebra, Volume 221 (2017) no. 5, pp. 1218-1241. \url{https://doi.org/10.1016/j.jpaa.2016.09.007}.
\bibitem[22]{published22} Fløystad, Gunnar; Orlich, Milo Triangulations of polygons and stacked simplicial complexes: separating their Stanley-Reisner ideals, J. Algebraic Combin., Volume 57 (2023) no. 3, pp. 659-686. \url{https://doi.org/10.1007/s10801-022-01174-7}.
\bibitem[23]{published23} Gangopadhyay, Rahul; Shannigrahi, Saswata Rectilinear crossings in complete balanced $d$-partite $d$-uniform hypergraphs, Graphs Combin., Volume 36 (2020) no. 3, pp. 905-911. \url{https://doi.org/10.1007/s00373-020-02163-y}.
\bibitem[24]{GaHiPe} Gasharov, Vesselin; Hibi, Takayuki; Peeva, Irena Resolutions of $\mathbf{a}$-stable ideals, J. Algebra, Volume 254 (2002) no. 2, pp. 375-394. \url{https://doi.org/10.1016/S0021-8693(02)00083-2}.
\bibitem[25]{published25} Hall, Huntington Tracy Counterexamples in discrete geometry, Ph. D. Thesis, University of California, Berkeley (2004).
\bibitem[26]{published26} Hao, Hanson; Luo, Yuyuan; Saint Rain, Sterling Polarizations of equigenerated strongly stable ideals, 2023 talk at Joint Mathematics Meetings (JMM 2023).
\bibitem[27]{published27} Hatcher, Allen Algebraic topology, Cambridge University Press, Cambridge, 2002, xii+544 pages.
\bibitem[28]{HeHi} Herzog, Jürgen; Hibi, Takayuki Monomial ideals, Graduate Texts in Mathematics, 260, Springer-Verlag London, Ltd., London, 2011, xvi+305 pages. \url{https://doi.org/10.1007/978-0-85729-106-6}.
\bibitem[29]{HT-cones} Herzog, Jürgen; Takayama, Yukihide Resolutions by mapping cones, Homology Homotopy Appl., Volume 4 (2002) no. 2, pp. 277-294. \url{https://doi.org/10.4310/hha.2002.v4.n2.a13}.
\bibitem[30]{published30} Huneke, Craig; Ulrich, Bernd The structure of linkage, Ann. of Math. (2), Volume 126 (1987) no. 2, pp. 277-334. \url{https://doi.org/10.2307/1971402}.
\bibitem[31]{published31} Jevtić, Filip D.; Timotijević, Marinko; Živaljević, Rade T. Polytopal Bier spheres and Kantorovich-Rubinstein polytopes of weighted cycles, Discrete Comput. Geom., Volume 65 (2021) no. 4, pp. 1275-1286. \url{https://doi.org/10.1007/s00454-019-00151-5}.
\bibitem[32]{published32} Kalai, Gil Many triangulated spheres, Discrete Comput. Geom., Volume 3 (1988) no. 1, pp. 1-14. \url{https://doi.org/10.1007/BF02187893}.
\bibitem[33]{published33} Lee, Carl W. Subdivisions and triangulations of polytopes, Handbook of discrete and computational geometry (CRC Press Ser. Discrete Math. Appl.), CRC, Boca Raton, FL, 1997, pp. 271-290.
\bibitem[34]{published34} Lu, Dancheng; Wang, Zexin The resolutions of generalized co-letterplace ideals and their powers, J. Algebra, Volume 673 (2025), pp. 321-350. \url{https://doi.org/10.1016/j.jalgebra.2025.02.029}.
\bibitem[35]{published35} Matoušek, Jiří Using the Borsuk-Ulam theorem: Lectures on topological methods in combinatorics and geometry, Universitext, Springer-Verlag, Berlin, 2003, xii+196 pages.
\bibitem[36]{published36} Migliore, Juan C. Introduction to liaison theory and deficiency modules, Progress in Mathematics, 165, Birkhäuser Boston, Inc., Boston, MA, 1998, xiv+215 pages. \url{https://doi.org/10.1007/978-1-4612-1794-7}.
\bibitem[37]{published37} Miller, Ezra Alexander duality for monomial ideals and their resolutions, 1998. \url{https://arxiv.org/abs/math/9812095}.
\bibitem[38]{published38} Miller, Ezra The Alexander duality functors and local duality with monomial support, J. Algebra, Volume 231 (2000) no. 1, pp. 180-234. \url{https://doi.org/10.1006/jabr.2000.8359}.
\bibitem[39]{MiSt} Miller, Ezra; Sturmfels, Bernd Combinatorial commutative algebra, Graduate Texts in Mathematics, 227, Springer-Verlag, New York, 2005, xiv+417 pages.
\bibitem[40]{published40} Mohammadi, Fatemeh; Pascual-Ortigosa, Patricia; Sáenz-de-Cabezón, Eduardo; Wynn, Henry P. Polarization and depolarization of monomial ideals with application to multi-state system reliability, J. Algebraic Combin., Volume 51 (2020) no. 4, pp. 617-639. \url{https://doi.org/10.1007/s10801-019-00887-6}.
\bibitem[41]{MoKa} Moradi, Somayeh; Khosh-Ahang, Fahimeh On vertex decomposable simplicial complexes and their Alexander duals, Math. Scand., Volume 118 (2016) no. 1, pp. 43-56. \url{https://doi.org/10.7146/math.scand.a-23295}.
\bibitem[42]{Mun} Munkres, James R. Topology, Prentice Hall, Inc., Upper Saddle River, NJ, 2000, xvi+537 pages.
\bibitem[43]{published43} Murai, Satoshi Spheres arising from multicomplexes, J. Combin. Theory Ser. A, Volume 118 (2011) no. 8, pp. 2167-2184. \url{https://doi.org/10.1016/j.jcta.2011.04.015}.
\bibitem[44]{published44} Nagel, Uwe; Reiner, Victor Betti numbers of monomial ideals and shifted skew shapes, Electron. J. Combin., Volume 16 (2009) no. 2, Paper no. 3, 59 pages. \url{https://doi.org/10.37236/69}.
\bibitem[45]{published45} Nematbakhsh, Amin Linear strands of edge ideals of multipartite uniform clutters, J. Pure Appl. Algebra, Volume 225 (2021) no. 10, Paper no. 106690, 21 pages. \url{https://doi.org/10.1016/j.jpaa.2021.106690}.
\bibitem[46]{Olt-Const} Olteanu, Anda Constructible ideals, Comm. Algebra, Volume 37 (2009) no. 5, pp. 1656-1669. \url{https://doi.org/10.1080/00927870802210001}.
\bibitem[47]{published47} Ovchinnikov, Sergei Graphs and cubes, Universitext, Springer, New York, 2011, xiv+287 pages. \url{https://doi.org/10.1007/978-1-4614-0797-3}.
\bibitem[48]{published48} Peskine, C.; Szpiro, L. Liaison des variétés algébriques. I, Invent. Math., Volume 26 (1974), pp. 271-302. \url{https://doi.org/10.1007/BF01425554}.
\bibitem[49]{published49} Postnikov, Alexander Permutohedra, associahedra, and beyond, Int. Math. Res. Not. IMRN (2009) no. 6, pp. 1026-1106. \url{https://doi.org/10.1093/imrn/rnn153}.
\bibitem[50]{published50} Sturmfels, Bernd Gröbner bases and convex polytopes, University Lecture Series, 8, American Mathematical Society, Providence, RI, 1996, xii+162 pages. \url{https://doi.org/10.1090/ulect/008}.
\bibitem[51]{published51} Yanagawa, Kohji Alexander duality for Stanley-Reisner rings and squarefree $\mathbf{N}^n$-graded modules, J. Algebra, Volume 225 (2000) no. 2, pp. 630-645. \url{https://doi.org/10.1006/jabr.1999.8130}.
\end{thebibliography}

\end{document}
